\documentclass[12pt,oneside]{article}

\usepackage[margin=1in]{geometry}
\usepackage{times}
\usepackage[hidelinks]{hyperref}

\pagestyle{empty}

\begin{document}

\begin{flushleft}

MISJustice Alliance \\
On behalf of Elvis Ryland Nuno and other victims \\
\today

\vspace{0.3in}

Federal Bureau of Investigation \\
Billings Field Office \\
200 North 3rd Avenue North \\
Billings, MT 59101 \\
Attn: Montana Resident Agencies (Including Missoula RA)

\vspace{0.3in}

\textbf{RE: Comprehensive Criminal Investigation Referral --- Montana-Based Conduct} \\
\textbf{Pattern of Institutional Corruption, RICO Violations, Civil Rights Conspiracy} \\
\textbf{YWCA Missoula, Missoula Police Department, Missoula County Attorney, Associated Officials} \\
\textbf{Federal Jurisdiction: 18 U.S.C. § 1962 (RICO); 18 U.S.C. § 241, 242; 42 U.S.C. § 1983, 1985} \\
\textbf{Coordination with FBI Seattle, USAO District of Montana, USAO Western District of Washington}

\vspace{0.3in}

Dear Special Agent in Charge:

\vspace{0.2in}

MISJustice Alliance respectfully submits this \textbf{comprehensive criminal investigation referral} to the FBI Billings Field Office regarding a decade-long pattern of institutional corruption, RICO violations, and systematic civil rights abuses centered in Missoula, Montana. This Montana-focused cover letter accompanies and complements a parallel referral sent to the FBI Seattle Field Office Civil Rights squad and is intended to ensure that all Montana-based conduct falls squarely under the jurisdiction of the newly established Billings Field Office and its resident agencies.

\section*{SCOPE OF MONTANA-BASED INVESTIGATION REQUEST}

While the broader submission documents an interstate pattern spanning Washington State and Montana (2015--2025), this letter highlights conduct occurring within Montana that warrants investigation by FBI Billings and the Missoula Resident Agency. The requested investigation targets:

\begin{enumerate}
\item \textbf{Institutional Enterprise in Montana:} YWCA Missoula and its governance structure.
\item \textbf{Law Enforcement Actors:} Missoula Police Department officers and supervisors, including Detective Connie Brueckner and Officer Ethan Smith.
\item \textbf{Prosecutorial Actors:} Missoula County Attorney's Office personnel who relied on demonstrably false narratives and suppressed exculpatory evidence.
\item \textbf{Individual Perpetrators:} Danielle Christine Chard and E'Lise Marie Chard, to the extent their conduct occurred or was executed through Montana institutions.
\item \textbf{Pattern Investigation:} Expanded inquiry into institutional corruption in Missoula and related Montana jurisdictions affecting multiple victims beyond Elvis Ryland Nuno.
\end{enumerate}

\section*{FEDERAL JURISDICTION (MONTANA COMPONENT)}

\subsection*{RICO Violations (18 U.S.C. § 1962)}

\textbf{Enterprise:} YWCA Missoula (nonprofit corporation), Missoula Police Department, and Missoula County Attorney's Office operating together as an association-in-fact enterprise.

\textbf{Pattern of Racketeering Activity:} 20+ predicate acts documented across 2012--2025, including:

\begin{itemize}
\item \textbf{Mail/Wire Fraud (18 U.S.C. § 1341, § 1343):} Federal and state grant applications and sponsor reports containing material misrepresentations regarding services, outcomes, staffing, and client protections.
\item \textbf{Obstruction of Justice (18 U.S.C. § 1503, § 1512):} Suppression of complaints, retaliation against whistleblowers and complainants, and interference with reporting to law enforcement and regulators.
\item \textbf{HIPAA Criminal Violations (42 U.S.C. § 1320d-6):} Systematic disclosure of protected health information by YWCA Missoula staff, including sensitive counseling information shared with abusers.
\item \textbf{Witness Tampering (18 U.S.C. § 1512):} Intimidation and retaliation against individuals who complained about YWCA Missoula or reported misconduct by affiliated officers and prosecutors.
\end{itemize}

\textbf{Interstate Commerce Impact:} Federal grant funds, interstate communications with Washington-based co-conspirators, and use of electronic platforms to coordinate false reports and case strategies.

\subsection*{Civil Rights Violations (18 U.S.C. § 241, § 242; 42 U.S.C. § 1983, § 1985)}

\begin{itemize}
\item \textbf{Fourth Amendment:} Warrantless home entries, unreasonable searches and seizures, and false imprisonment directed at Mr. Nuno and others in Missoula.
\item \textbf{First Amendment:} Retaliation following Mr. Nuno's formal YWCA complaint and subsequent criticism of institutional misconduct.
\item \textbf{Fourteenth Amendment:} Malicious prosecution, denial of due process, and discriminatory treatment in charging and bail decisions.
\item \textbf{Civil Rights Conspiracy:} Coordinated actions by YWCA Missoula staff, MPD officers, and local prosecutors to deprive individuals of equal protection and to silence critics.
\end{itemize}

\section*{EVIDENCE BEYOND NUNO CASE: ENDEMIC CORRUPTION IN MONTANA}

The enclosed \textbf{YWCA Institutional Corruption Supplemental} (Document \#11) establishes RICO predicate acts \textbf{independent of the Nuno case}, demonstrating a broader criminal enterprise rooted in Montana:

\subsection*{Multiple Documented Victims in Montana}

\begin{itemize}
\item \textbf{Arthur Brown:} YWCA Missoula submitted defamatory letters to a Montana court recommending permanent restraining orders without adequate factual basis, resulting in separation of a father from his children through fraud upon the court.
\item \textbf{Anonymous Client \#1:} HIPAA violations and ADA discrimination stemming from YWCA Missoula's disclosure of confidential counseling information to an abusive ex-partner through staff personal connections.
\item \textbf{Anonymous Client \#2:} Discriminatory service provision and federal grant fraud, where outcomes reported to sponsors were inconsistent with actual treatment of families.
\item \textbf{Multiple Employees:} 2023--2024 reviews describing racist undertones, untrained staff, and dangerous conditions for vulnerable clients, all within YWCA Missoula operations.
\end{itemize}

\subsection*{Institutional Governance Capture (Missoula)}

\textbf{Detective Connie Brueckner} simultaneously served as:

\begin{itemize}
\item A YWCA Missoula Board Member at-large.
\item The lead Missoula Police Department investigator of complaints against YWCA employees.
\item A primary actor in unconstitutional actions against complainants, including Mr. Nuno.
\end{itemize}

This conflict of interest prevented impartial review of complaints, concealed Brady material from criminal defendants, and effectively insulated YWCA Missoula and MPD actors from accountability within Montana.

\subsection*{Federal Oversight Context in Montana}

\begin{itemize}
\item \textbf{2012--2016 DOJ Investigation:} Missoula was identified as having systemic failures in sexual-assault response, and YWCA Missoula was designated as a central “community partner” in reform, even as governance conflicts and civil rights violations persisted.
\item \textbf{2023 LifeGuard Group Matter:} A Montana DOJ Agent documented “concerning” strip search allegations involving a survivor and child at a facility tied to local law enforcement leadership and significant philanthropic funding; no criminal investigation followed, exemplifying broader oversight failures.
\end{itemize}

\subsection*{Institutional Refusal to Respond to Formal Notices}

MISJustice Alliance sent formal legal notices to YWCA Missoula on \textbf{October 12, 2025} and \textbf{December 5, 2025}, detailing conflicts, retaliation, confidential-information misuse, and grant fraud. The December notice was copied to senior YWCA staff, the Missoula Mayor, Missoula Police Chief, County Attorney, and other officials. \textbf{To date, YWCA Missoula has offered no response or acknowledgment}, reinforcing evidence of deliberate indifference and institutional capture that cannot be remedied through local channels.

\section*{RELATIONSHIP TO SEATTLE FBI AND INTERSTATE PATTERN}

This Montana-focused referral is part of an integrated, dual-jurisdiction pattern:

\begin{itemize}
\item Washington charges and actions (Seattle/Edmonds/King County) were influenced by and interwoven with Montana-based reports, narratives, and decisions.
\item Montana prosecutions and police actions relied on Washington matters to justify ongoing scrutiny and harassment of Mr. Nuno.
\item The FBI Seattle Field Office Civil Rights squad is being asked to lead on interstate coordination and Washington-based actors, while the FBI Billings Field Office is being asked to lead on Montana-based conduct and resident-agency operations.
\end{itemize}

Parallel referrals have been made to the \textbf{U.S. Attorney's Office for the District of Montana}, the \textbf{U.S. Attorney's Office for the Western District of Washington}, and the \textbf{DOJ Criminal Division Public Integrity Section} for potential charges against public officials and prosecutors.

\section*{REQUESTED FEDERAL ACTIONS (MONTANA FOCUS)}

\begin{enumerate}
\item \textbf{Opening of a Montana-Based Investigation:} Initiate an FBI Billings/Missoula RA investigation into YWCA Missoula, Missoula Police Department, and Missoula County Attorney's Office conduct affecting Mr. Nuno and other victims.
\item \textbf{Coordination with FBI Seattle:} Coordinate with the Seattle Field Office Civil Rights squad regarding the interstate pattern and Washington-based co-conspirators.
\item \textbf{Coordination with USAO District of Montana:} Work with the U.S. Attorney to assess RICO, civil rights, and related charges and to consider grand jury presentation in the District of Montana.
\item \textbf{Referral to DOJ Public Integrity Section:} Confer with PIN concerning potential public corruption and prosecutorial misconduct by Montana officials.
\item \textbf{Victim and Witness Protection:} Assess and address retaliation risks for Mr. Nuno and other reporting victims who may provide testimony or additional documentation.
\end{enumerate}

\section*{SUBMISSION CONTENTS (11 Documents)}

This cover letter is part of the same unified packet sent to all listed federal and state authorities and includes:

\begin{enumerate}
\item Case Summary.
\item FBI Seattle Field Office Civil Rights Cover Letter.
\item This FBI Billings Field Office Cover Letter (Montana-focused).
\item Montana Attorney General Cover Letter.
\item Washington Attorney General Cover Letter.
\item U.S. Attorney, District of Montana Cover Letter.
\item U.S. Attorney, Western District of Washington Cover Letter.
\item DOJ Criminal Division Public Integrity Section Cover Letter.
\item Danielle Chard Criminal Report.
\item E'Lise Chard Criminal Report
\item YWCA Institutional Corruption Supplemental; Elvis Nuno Sworn Declaration; and comprehensive evidentiary documentation (police reports, court records, correspondence, and financial records).
\end{enumerate}

\section*{EVIDENTIARY STANDARD AND URGENCY}

This packet reflects \textbf{thousands of hours} of collaborative work by MISJustice Alliance members, including extensive deep records research and interviews with multiple victims. The documentation includes sworn declarations, official records, corroborating witness statements, and financial analyses. The evidence is comprehensive and ready for immediate grand jury presentation in the District of Montana, with particular urgency for predicate acts nearing statute-of-limitations thresholds. The evidence is comprehensive and ready for immediate grand jury presentation.

We respectfully request an initial response via digital communication outlining appropriate next steps, sent to \texttt{advocacy@misjusticealliance.org}, so that investigation planning, grand jury coordination, and inter-office cooperation can proceed while protecting the anonymity and safety of MISJustice Alliance members and reporting victims.

\vspace{0.3in}

\noindent
Respectfully submitted,

\vspace{0.3in}

\noindent
MISJustice Alliance \\
On behalf of Elvis Ryland Nuno and other victims

\vspace{0.2in}

\noindent
Contact: [\textit{Contact information will be provided only through MISJustice Alliance after confirmation of concrete plan to prevent any retaliation or adverse action against reporting victims and witnesses.}]

\end{flushleft}

\end{document}
